Зафиксируем используемые понятия индекса. Если $v$ --- непустое
конечное слово и $s\ge1$ --- рациональное число с $s|v|\in\mathbb N$,
то $v^s$ означает префикс слова $v^\omega$ длины $s|v|$.
Для бесконечного слова $u$ положим
\[
 \operatorname{ind}_u(v)
 =\sup\{s\in\mathbb Q:\ s\ge1,\ s|v|\in\mathbb N,\ v^s\in\mathcal L_u\},
 \]
 \[
 E(u)=\sup_{v\in\mathcal L_u\setminus\{\varepsilon\}}\operatorname{ind}_u(v).
\]
Асимптотический критический показатель определяется как
\[
 E^*(u)=\lim_{M\to\infty}
 \sup_{\substack{v\in\mathcal L_u\\ |v|\ge M}}\operatorname{ind}_u(v).
\]
Предел существует в $[1,+\infty]$, поскольку выражение под
пределом не возрастает по $M$. В дальнейшем запись
$\operatorname{ind}(u)$ означает $E(u)$, а
$\operatorname{ind}^*(u)$ --- именно $E^*(u)$.

Перестановка букв сохраняет показатели $E$ и $E^*$.
Поэтому для штурмова слова с угловым коэффициентом $\alpha$
можно нормировать алфавит так, чтобы используемый коэффициент
был $\beta=\min\{\alpha,1-\alpha\}<1/2$.
В следующем утверждении $\varphi=(1+\sqrt5)/2$.

\begin{theorem}[Классические формулы для штурмовых слов
{\cite{DamanikLenz2002,Peltomaki2015,AdamczewskiAllouche2007}}]\label{thmconsn1}
Пусть $u$ --- штурмово слово с угловым коэффициентом
$\alpha\in(0,1)\setminus\mathbb Q$.
Положим $\beta=\min\{\alpha,1-\alpha\}=[0;a_1,a_2,\ldots]$,
$q_{-1}=0$, $q_0=1$ и
$q_{n+1}=a_{n+1}q_n+q_{n-1}$.
Тогда
\begin{align}
 E(u)&=\max\left\{a_1,\
       \sup_{n\ge1}\left(2+a_{n+1}+\frac{q_{n-1}-2}{q_n}\right)\right\},
       \label{eq:correct-index}\\
 E^*(u)&=2+\limsup_{n\to\infty}\frac{q_{n+1}}{q_n}
       =2+\limsup_{n\to\infty}[a_n;a_{n-1},\ldots,a_1].
       \label{eq:correct-asymptotic-index}
\end{align}
В частности, $E^*(u)\ge2+\varphi$.
\end{theorem}
\begin{proof}
Поскольку обмен букв не меняет показатели, работаем с нормированным
языком углового коэффициента $\beta<1/2$.
Используем следующую классическую структурную информацию о степенях
в штурмовом языке
\cite[теорема~4.5, следствие~4.6 и доказательство теоремы~4.7]{Peltomaki2015}:
\begin{itemize}
 \item если примитивное слово $v$ (не являющееся нетривиальной
 степенью более короткого слова) удовлетворяет
 $v^3\in\mathcal L_u$, то либо $v=0$, либо $|v|=q_n$
 для некоторого $n\ge1$;
 \item $\operatorname{ind}_u(0)=a_1$, а наибольший дробный показатель
 среди примитивных факторов длины $q_n$ равен
 \[
 e_n=2+a_{n+1}+\frac{q_{n-1}-2}{q_n}
     =2+\frac{q_{n+1}}{q_n}-\frac2{q_n}.
 \]
 Для каждого $n\ge1$ соответствующий фактор и его степень
 действительно встречаются.
\end{itemize}
Это используемый внешний комбинаторный результат; последующий вывод
асимптотической формулы приводится полностью.

Сначала заметим, что
$L=\limsup_{n\to\infty}q_{n+1}/q_n\ge\varphi$.
Если бы $L<\varphi$, можно было бы выбрать $c$ с
$L<c<\varphi$. Для всех достаточно больших $n$
$q_{n+1}/q_n\le c$, но рекурсия тогда даёт
\[
 \frac{q_{n+2}}{q_{n+1}}
 =a_{n+2}+\frac{q_n}{q_{n+1}}
 \ge1+\frac1c>c,
\]
что противоречит выбору $c$. Поэтому, поскольку $q_n\to\infty$,
\[
 \limsup_{n\to\infty}e_n=2+L\ge2+\varphi>3.
\]

Каждое непустое слово $w$ единственным образом записывается
$w=v^t$, где $v$ примитивно, а $t\ge1$ --- целое.
Из определения степеней следует
$\operatorname{ind}_u(w)\le\operatorname{ind}_u(v)/t$.
Если $\operatorname{ind}_u(w)>3$, то $v^3$ встречается.
По структурному результату $v=0$ или $|v|=q_n$;
показатель не превосходит соответственно $a_1$ или $e_n$.
Показатели $\le3$ не могут увеличить супремум, поскольку
$\sup e_n>3$. Обратные оценки обеспечиваются существованием
$0^{a_1}$ и степеней с показателями $e_n$.
Это доказывает~\eqref{eq:correct-index}.

Для асимптотической формулы существование степеней с корнями
длины $q_n\to\infty$ даёт
$E^*(u)\ge\limsup e_n=2+L$.
Если $L=+\infty$, равенство уже доказано.
В противном случае рассмотрим любую последовательность факторов
$w_j$ с $|w_j|\to\infty$ и показателями, стремящимися к $E^*(u)$.
Достаточно рассматривать показатели $>3$.
Для примитивного корня $v_j$ слова $w_j$ тогда имеем
либо $v_j=0$, либо $|v_j|=q_{n_j}$.
Первый случай не может встречаться при неограниченных $|w_j|$:
длина постоянного фактора не превосходит $a_1$.
Если индексы $n_j$ ограничены, то для каждого такого индекса
показатель примитивного корня конечен, и условие
$\operatorname{ind}_u(w_j)>3$ ограничивает также
$t_j=|w_j|/|v_j|$ величиной $e_{n_j}/3$.
Это вновь противоречит $|w_j|\to\infty$.
Следовательно, $n_j\to\infty$ после выбора подпоследовательности, и
\[
 \limsup_j\operatorname{ind}_u(w_j)
 \le\limsup_j e_{n_j}\le2+L.
\]
Таким образом, $E^*(u)=2+L$.
Наконец, рекурсия знаменателей даёт зеркальную формулу
$q_{n+1}/q_n=[a_{n+1};a_n,\ldots,a_1]$;
сдвиг индекса верхнего предела доказывает последнее равенство
в~\eqref{eq:correct-asymptotic-index}.
\end{proof}

\noindent Для слова Фибоначчи
$\beta=\varphi^{-2}=[0;2,1,1,\ldots]$ и $q_n=F_{n+2}$.
Из формулы Бине
$F_{n+3}-\varphi F_{n+2}=\psi^{n+2}$, где $\psi=-\varphi^{-1}$,
следует
\[
 e_n=2+\varphi+\frac{\psi^{n+2}-2}{F_{n+2}}<2+\varphi,
 \qquad e_n\longrightarrow2+\varphi.
\]
Поэтому $E(\mathbf f)=E^*(\mathbf f)=2+\varphi$.
Показатель $1$ не является возможным значением $E^*$ штурмова слова.
